\begin{textsample}{NEG Dim 3 – mt – Score -3.00 – mt/mt\_em\_37.txt}  \label{ex:f3_neg_206}
7.4 Realidades cotidianas enfrentadas pelas \textbf{Juventudes} do Ensino Médio O grande desafio a ser enfrentado pelas \textbf{juventudes} da Etapa Ensino Médio é a formação para o mundo do trabalho.
É importante destacar que essa concepção vê o trabalho como práxis humana e não como produtor de mercadorias.
A preparação para o mundo do trabalho deve estar articulada com conhecimentos científicos, culturais e sócio-históricos, no âmbito da emancipação humana.
Os diversos grupos de jovens possuem características e habilidades diferenciadas de acordo com seus saberes e experiências constituídos ao longo de suas vivências e sua condição socioeconômica e cultural.
Sendo assim, é necessário valorizar as diversidades culturais, étnicas e de gênero dos estudantes, respeitando suas especificidades e potencialidades.
Ao considerarmos as diversas realidades cotidianas enfrentadas pelos jovens na atualidade, o DRC-MT do Ensino Médio tem como objetivo a implementação desta etapa, partindo do princípio de habilidades e \textbf{competências} que atendam os anseios educativos das \textbf{juventudes}, respeitando o protagonismo estudantil e seu projeto de vida.
No entanto, para que isso se efetive, é indispensável a melhoria das condições de acesso e permanência dos estudantes, com maiores investimentos em equipamentos, ampliação dos espaços físicos, investimentos em materiais pedagógicos e formação continuada dos profissionais da educação : professores, técnicos administrativos e apoios administrativos educacionais.
O Ensino Médio só será efetivamente democrático quando os Projetos Políticos Pedagógicos forem eficientes e eficazes para as \textbf{juventudes}, independentemente de sua classe econômica, cultural e social.
Os menos favorecidos devem ter acesso à cultura e a conhecimentos científicos através do desenvolvimento de \textbf{competências} que são inerentes à escolarização, portanto, as metodologias e aprendizagens devem ser variadas, significativas e adequadas aos Projetos Político Pedagógicos das escolas.
Deste modo, deve -se atender aos anseios da periferia, classe alta, da educação escolar indígena e quilombola, das pessoas com deficiência e dos sujeitos do campo.
Só assim, teremos uma escola justa democrática e inclusiva, em que todos estejam em condições de identificar, compreender e suprir, ao longo de suas vidas, suas necessidades educacionais.

% matched lemmas: competência, juventude
\end{textsample}
